\section{Introduction}

% The goal is to not write a passive "policy paper" like the UN, but to utilize EO data and climate modelling as means to guide the poltically informed "conversations" in the dissertation.

\subsection{The Urban Climate}

% Synthesis and summary of social science and science based literature of the urban climate
    % scientific explanations that give way to social-historical results
In the start of \textit{Urban Climates} by urban climatologists \textcite{urban_climates}, the city was described as a "wide array of surface elements of very different sizes and compositions, [with] the vast majority of them [] alien to the natural landscape and include pulses of energy, water, gases and particles \textbf{controlled by people} rather than geophysical activity"(p. xix, author emphasis). This is the crux of understanding the urban climate, whereby the built environment is inherently artificial, in which it's processes are clearly anthropogenic, and where studies must bridge this gap, recognizing that research of the urban environment is thus primarily a study of human processes first and foremost. The traditional division between the "natural" and "social" sciences, and within geography as "physical" and "human" falls apart in the city. This dissertation serves as an exploration of this interdisciplinarity, critically utilizing climate modelling as means of understanding the geospatial inequalities of the \Gls{uhi} in Greater Kuala Lumpur.  %/cite. Page 55 un_2018_wup, 60.4% urbanpop, 27.8% in city >1mill

The world is facing a multitude of crises, environmental, social and economic, that are especially embodied in the urban city. The \textcite{un_2018_wup} estimates that 60\% of the world population, consisting of 5.2 billion people, will be living in urban areas by 2030 with nearly half living in metropolitan areas of greater than a million people. The Global South in particular is facing unprecedented and rapid urbanization, whereby cities like Karachi, Jakarta and Dar es Salaam will overtake well-known Northern agglomerations like New York, Paris and London, with more than 100 settlements having a population of greater than 5 million by 2030. These urban agglomerations have grown to such an extent to have noticeable effects on their regional climates, with numerous studies having established that urban environments are prone to rising surface and urban temperatures, alterations of precipitation patterns and intensity at both a sub-daily and seasonal scales, the destruction of natural ecosystems, and a complicated interaction of this large urban ecosystem with it's peri-urban 'semi-natural' surroundings. 

Understanding the biophysical and social processes that dictate and govern the living conditions for a majority of mankind becomes not only a matter of scientific curiosity but a social necessity to equip those most affected with the knowledge and tools to rewrite their own circumstances (\cite{crit_physicalgeo}). Echoing this sentiment, there is an urgent need to carry forward the goal of 'common but differentiated responsibilities and respective capabilities' in response to anthropogenic climate change (\cite{paris_agreement}). However, in practical terms, there is still a consistent "data gap" between the Global North and South, which means there is considerable difficulty for research to be equitable and consistently accurate between different regions. In addition to this data gap, Southern nations face additional challenges in adapting and responding effectively to the rapid growth of their cities, which often coincide with notions of economic progress and development (\cite{urban_resilience}). Socioeconomic and political factors continuously limit the ability for governments and other institutions to setup and gather data in many of these "informal" urban growths, and scholarship bias to Northern urban environments leads to a lack of established literature and relevant procedures to design comparable and reproducible studies of Southern urban climates.


The call for a critical physical geography by \textcite{crit_physicalgeo} is especially relevant for urban climate research, where the resulting environmental conditions and research itself is a direct result of political-economic decision-making. Initiatives for utilizing academically rigorous, neutral and globally available datasets, with open-source and replicable methodology, is prudent in bridging this gap while enabling accessible 'citizen science'.

This dissertation utilizes an open-source and standardized methodology of simulating the the urban climate of Greater Kuala Lumpur to further inform discussions regarding the unequal geospatial effects of the \Gls{uhi} and their potential social effects in the Greater Kuala Lumpur urban form and function. 



% expounding general example of city environments and their impacts on climactic conditions
 
% sociopolitical risks, urban planning,



\subsection{Urban Climate Modelling}
% Urban Climate Science? The field of Urban Climate?
% The scientific (ie. conceptual) basis of Urban Climate

In the era of big data, and a proliferation of Earth System Science modelling through global \Gls{eo} data, this full head-on engagement with derived information without on the ground measurements  

There has been a lack of synthesizing \Gls{eo} datasets in informing decision-making at a sub-national level due to the complexity of downsizing the data to a smaller level capable of accurately portraying 

% biophysical factors affecting the urban climate

% rise of computational modelling, development of procedures and related modelling work
Climate modelling over the past couple of decades have seen a revolutionary breakthrough as advanced in technology has made it possible to calculate 

Land Surface Models aim to understand the interactions between the surface and atmospheric layers.
% what is it? what are it's mechanisms? what makes it different?
UBL, UCL

%the role of the urban form

% the meso and micro scale

\subsection{Health impacts of Heat Stress}



\subsection{Current Urban Heat research}
The IPCC AR6 report has confirmed that heat stress remains one of the many challenges faced by growing urban environments

The urban heat island has been an established concept in urban climatology for more than half a century (\cite{oke_1973}), primarily due to observed experience and the growth of urban environments. One key way to conceptualize urban heat is through the energy balance at the surface, in which urban environments, having morphometric properties different to that of rural environments, lead to an increase in energy intake and thus air temperature in a city (\cite{urban_energy_balance}). Most research has focused on Northern cities and so understanding the mechanisms of urban heat are still scant in tropical climate regimes where intra-annual temperatures vary less than diurnal ranges. Furthermore, tropical climate regimes face seasonal variations in precipitation which will affect the mentioned urban energy balance equation by introducing varying levels of latent heat flux.

A systemic literature review carried out by \textcite{kim_2021} has found that most urban heat studies can be categorized into two main methodologies, namely experimental or empirical. Due to the difficulty of getting accurate and widespread meteorological data for many of the sites concerned, many research papers that utilize experimental designs
The \Gls{uhi} phenomenon has been noted since xy 

Many papers utilize numerical simultions such as that of the WRF model

As many urban climatologists have noted, \Gls{uhi} suffers from a lack of standardization of categories and definitions. In this paper, the focus will be on the Urban Heat Island at 2m air temperature, following the conventional standards for urban heat island studies,



The dissertation will fall under the experimental study, with specific focus on se

(insert UHII equation)
Urban Heat Island Intensity 



% urban heat island effect, and associated literature

% definitions
% Use standardized or proposed standardized UHI units and symbols!!! 

% understanding the differences between UHI from surface studies ie. radiative measurements, UHI in the UCL (10-20m dependning on local morpholgoy), UHI in the UBL (<100m). THe study focuses on T2M,

% Ditstinction between surface UHI and etc

% health effects

% Limitations moved to under current heat reserach

A lot of climate modelling not only uses globl big data networks of Earth Observation data, this leads to layers of abstraction t

% local scale (1km and sub 500m) climate modelling, lack of data and high computer usage



\newpage
\section{Site of Study}
    
\subsection{Social and historical background}

\gls{GKL}, also known as Klang Valley (\cite{epu}), is a large urban conurbation in Malaysia consisting of multiple local authorities including that of the capital city of it's name, the Federal Territory of Kuala Lumpur. The city is currently estimated to have a population of nearly 9million, with future projections 


population

brief history, fishing vilalge etc
    % urbanization, expansion
    % urban planning strategies, neoliberal private-led development
    % high rises 

\subsection{Climactic and environmental background}

\Gls{gkl} is located on the confluence of the Klang and Gombak rivers, with an extensive urban sprawl that reaches the coastline at Port Klang. As a result, while not directly at the coast, the city experiences notable coastal influences in it's climate, especially that of land-sea breezes. Additionaly, the Titiawangsa mountain range serves as a great central divider of Peninsular Malaysia, located to the city's northeast. 

%add topograhical map of peninsular Malaysia, with state boundaries and box of site of study
%zoom in figure 

\Gls{gkl} falls under the meteorological \Gls{mc}, in which diurnal patterns remain the focal point and main driver of larger scale climactic processes. 

% add satelite and topographical image of the study area
    % overlay with model area?
% climactic and meteorological background of the Greater Kuala Lumpur
    % the two monsoon seasons
    % rainfall patterns

% Exisitng urban heat studies
% \subsubsection

In the context of Malaysia, most studies still employ global datasets with sparse data resolution without modelled downscaling . Gl

\subsection{Aims and Objectives}

{\textit\large{Main research question}}
- How does local-scale urban morphology affect the spatio-temporal intensity of the urban heat islands in the tropical climate of Greater Kuala Lumpur?

{\textit\large{Aims}}
- To learn and utilize Python for data science
- 
   % - *surface cover fractions: normal = through LCZ, modified = alteration of default cover fraction conversion parameter
   % - *Independent variable: GKL land cover over time*
   % - *Dependent variable: near surface temperature (2m), relative humidity*


- What are the spatial variations of \Gls{uhii}?
    - *Which parts of GKL exhibit largest UHII? Does this pattern hold over time?*

- What are the diurnal and seasonal patterns of the UHII?



